\section{Dynamics, first feature: apo-MD RMSF}\label{sec:rmsf}

The B-factor test (\S\ref{sec:baseline}) was the static teaser; this is the measurement it was
standing in for. With the apo production runs landed, we replace crystallographic mobility with
its lattice- and antibody-free dynamic analogue: per-residue \textbf{root-mean-square fluctuation}
(RMSF) from the apo trajectories. If the weak rigidity lean of the B-factors ($\rho=-0.23$) was a
crystallographic artefact --- resolution, refinement, packing --- proper dynamics should wash it out;
if it was real, RMSF should sharpen it. Either way this is the first feature of the dynamics program,
and the one most directly comparable to a static baseline.

\paragraph{Setup.} Per-residue C$\alpha$ RMSF over the equilibrated portion of each $100$\,ns apo
trajectory (first $10$\,ns discarded; least-squares fit to the reference, so global tumbling is removed
before fluctuations are measured), z-normalised per structure to match the B-factor convention
(\texttt{analysis/rmsf\_dynamics\_baseline.py}; GROMACS \texttt{2023.2} on Gemini). Labels are kept
at their native granularity --- one column per (antigen, antibody) \emph{epitope}, never averaged over
antibodies, so lysozyme's three antibodies (\texttt{HyHEL-63}, \texttt{HyHEL-10}, and the opposite-face
\texttt{D1.3}) and VEGF's two are scored separately. Correlation is reported per epitope, pooled over
all scanned residues, and \textbf{leave-one-antigen-out} (LOAO): the honest generalisation number,
since residues within an antigen share a fold and a trajectory and our effective $N$ is the number of
antigens, not residues (\S\ref{sec:design}). The split is by antigen, not antibody --- lysozyme's three
antibodies share one apo trajectory, so an antibody-level split would leak features across train and test.

\paragraph{Which systems.} Six of the finished apo runs align cleanly to the labels
(\texttt{1AKI}, \texttt{2JEL}, \texttt{1BJ1}, \texttt{1JRH}, \texttt{1AHW}, \texttt{1HGU}); the MD
structure is numbered differently from the antibody-complex crystals the scans come from, so alignment
is \emph{verified by residue identity} --- each scanned wild-type residue must match the amino acid
actually present at that position in the simulated structure, not assumed from an offset. MT-SP1
(\texttt{1OKE}) fails this check (its protease numbering with activation-cleaved insertion codes does not
map to the \texttt{3BN9}/\texttt{3NPS} scan numbering) and is \textbf{excluded here}, flagged
[verify:\ numbering] rather than force-joined; the influenza HA run (\texttt{1HGG}) is the largest system
and had not finished at the time of this iteration. That leaves $9$ antibody-epitopes over $6$ antigens,
$97$ scanned residues.

\begin{figure}[H]\centering
\figorbox{rmsf_baseline.png}{0.98}
\caption{\textbf{Apo-MD flexibility as an epitope label.} \emph{Left:} Spearman $\rho$ between
per-residue apo RMSF and alanine-scan $\ddG$ for each antibody-epitope (coloured by antigen; dashed
line is the pooled value); the correlations scatter around zero with a slight negative lean.
\emph{Right:} all $97$ scanned residues, RMSF ($z$; right $=$ more mobile) vs.\ $\ddG$, colours as in
Fig.~\ref{fig:bfactor-baseline} --- the important (red) residues are spread across the mobility range,
not concentrated on the mobile side. An RMSF$\geq 0$ cutoff (the shaded band) flags $55$ of $97$ scanned
residues and catches $25$ of $48$ hotspots --- no better than a coin toss on this set.
\texttt{analysis/rmsf\_dynamics\_baseline.py}. MT-SP1 (\texttt{1OKE}) excluded (numbering); HA
(\texttt{1HGG}) pending.}\label{fig:rmsf-baseline}
\end{figure}

\paragraph{Result.} Proper dynamics does not rescue the flexibility hypothesis --- it confirms the
B-factor verdict on cleaner data. The pooled scanned Spearman between RMSF and $\ddG$ is
$\rho = -0.12$ ($p=0.24$, $n=97$), and the leave-one-antigen-out mean is $\rho = -0.15$ (over the five
antigens with $\geq 5$ scanned residues) --- statistically indistinguishable from zero, and the same
weak-rigid \emph{sign} the apo B-factors gave ($-0.23$). Per epitope the correlations run from
$\rho=-0.80$ (hGH, lysozyme \texttt{HyHEL-10}, both $n=4$) through $-0.66$ (lysozyme \texttt{D1.3}) to
$+0.29$ (IFN-$\gamma$ receptor); the sign is not even consistent across the three antibodies of a single
antigen. As a classifier the picture is the same: thresholding at RMSF$\geq 0$ recovers roughly half the
hotspots at roughly half the scanned residues --- the diagonal of no information. Whatever distinguishes
a binding-energy hotspot from its neighbours, it is not how much that residue moves.

\paragraph{Reading.} This is a negative result on the feature, but a positive one for the thesis of
\S\ref{sec:baseline}: neither static nor dynamic \emph{mobility} --- B-factor or RMSF --- carries the
hotspot signal, which is exactly what the energetic-isolation view predicts. Motion is the wrong
observable; what should matter is a residue's \emph{coupling} to the rest of the protein, not its
amplitude of fluctuation~\cite{serapian2020}. RMSF measures the amplitude and finds nothing, cleanly
and on real trajectories. The same runs that produced these fluctuations carry the observables that
the coupling view actually points at --- intra-protein nonbonded interaction energy, the
correlation structure of the motions, and mode participation --- and those are the next features, scored
against this same baseline.

\paragraph{Caveats.} (i) Single $100$\,ns trajectory per antigen; RMSF of the most flexible loops is
not fully converged at this length, though the hotspot residues (ordered, low-RMSF) are the stable part
of the estimate. (ii) Two epitopes rest on only $n=4$ scanned residues (hGH, lysozyme \texttt{HyHEL-10}),
so their individual $\rho$ values are indicative, not reliable; the pooled and LOAO numbers are the ones
to trust. (iii) MT-SP1 and HA are absent (numbering; run pending) --- the pooled value will be re-computed
when both join. (iv) As with the static baselines, unscanned residues are treated as non-epitope.
